\section*{Zdrojové kódy implementácie}
\addcontentsline{toc}{section}{Zdrojové kódy implementácie}
\label{sec:zdrojove_kody_arduino}

\begin{lstlisting}[language=C++, caption={Hlavný zdrojový kód pre Arduino UNO}, label={lst:main_cpp_appendix}]
#include <Arduino.h>

// a1 -> snimac1
// a2 -> snimac2
// pwm5 -> vent
// pwm6 -> spirala

// Pin configuration
constexpr uint8_t PWM_OUT_VENT = 5;
constexpr uint8_t PWM_OUT_SPIR = 6;
constexpr uint8_t ANALOG_IN_1 = A1;
constexpr uint8_t ANALOG_IN_2 = A2;

// Experiment timing
constexpr uint32_t TS_US = 100000UL;           
constexpr float TS_WARNING_S = 0.11f;
constexpr float EXPERIMENT_DURATION_S = 30.0f;

// Experiment state
bool experimentRunning = false;

// Global experiment timing
uint32_t experimentStartUs = 0;
uint32_t lastControlTickUs = 0;

// HELPER FUNCTIONS ////////////////////////////////////////////////////////////////

// Cubic polynomial evaluation
float cubic(float x, float a, float b, float c, float d) {
    return a * x * x * x + b * x * x + c * x + d;
}

// Clamp float value into range
float clampFloat(float value, float minValue, float maxValue) {
    if (value < minValue) {
        return minValue;
    }
    if (value > maxValue) {
        return maxValue;
    }
    return value;
}

// Clamp integer value into range
int clampInt(int value, int minValue, int maxValue) {
    if (value < minValue) {
        return minValue;
    }
    if (value > maxValue) {
        return maxValue;
    }
    return value;
}

// SERIAL COMMANDS //////////////////////////////////////////////////////////////////////////

// Stop all outputs
void stopOutputs() {
    analogWrite(PWM_OUT_VENT, 0);
    analogWrite(PWM_OUT_SPIR, 0);
}

// Start experiment and reset timers
void startExperiment() {
    experimentStartUs = micros();
    lastControlTickUs = experimentStartUs;
    experimentRunning = true;

    Serial.println("t,Ts,snimac1");
}

// Handle serial commands
void handleSerialCommands() {
    if (Serial.available() <= 0) {
        return;
    }

    String command = Serial.readStringUntil('\n');
    command.trim();

    if (command == "START") {
        if (!experimentRunning) {
            startExperiment();
        }
    } else if (command == "STOP") {
        experimentRunning = false;
        stopOutputs();
        Serial.println("STOPPED");
    }
}

void setup() {
    pinMode(PWM_OUT_VENT, OUTPUT);
    pinMode(PWM_OUT_SPIR, OUTPUT);
    pinMode(ANALOG_IN_1, INPUT);
    pinMode(ANALOG_IN_2, INPUT);

    Serial.begin(115200);

    stopOutputs();

    delay(300);
    Serial.println("READY");
}

// INIT control variables ///////////////////////////////////////

float vent = 0.0f;
float spirala = 0.0f;

// MAIN CONTROL LOOP ////////////////////////////////////////////

void loop() {
    handleSerialCommands();

    if (!experimentRunning) {
        return;
    }

    uint32_t nowUs = micros();
    uint32_t elapsedSinceLastTickUs = nowUs - lastControlTickUs;

    // Loop executes only when real elapsed time reached Ts
    if (elapsedSinceLastTickUs < TS_US) {
        return;
    }

    float timeTickPeriod = static_cast<float>(elapsedSinceLastTickUs) / 1000000.0f;
    float globalTime = static_cast<float>(nowUs - experimentStartUs) / 1000000.0f;

    // Update tick timestamp only when control loop actually runs
    lastControlTickUs = nowUs;

    // OUTPUT control variables (vent, spirala) ////////////////////

    // Limit control variables to 0..10
    vent = clampFloat(vent, 0.0f, 10.0f);
    spirala = clampFloat(spirala, 0.0f, 10.0f);

    // Convert vent -> PWM5 count using cubic calibration
    float pwm5Float = cubic(
        vent,
        0.05938291f,
        -0.27988543f,
        24.20478979f,
        -0.52522611f
    );

    // Convert spirala -> PWM6 count using cubic calibration
    float pwm6Float = cubic(
        spirala,
        0.08732906f,
        -0.39322033f,
        22.8496962f,
        -0.61564759f
    );

    // Round to nearest integer and clamp to valid PWM range
    int pwm5 = clampInt(static_cast<int>(pwm5Float + 0.5f), 0, 255);
    int pwm6 = clampInt(static_cast<int>(pwm6Float + 0.5f), 0, 255);

    // Output PWM
    analogWrite(PWM_OUT_VENT, pwm5);
    analogWrite(PWM_OUT_SPIR, pwm6);

    // INPUT processing //////////////////////////////////////////////

    // Read raw ADC values
    int a1 = analogRead(ANALOG_IN_1);
    int a2 = analogRead(ANALOG_IN_2);

    // Convert ADC -> snimac values using cubic calibration
    float snimac1 = cubic(
        static_cast<float>(a1),
        4.64044633e-10f,
        -8.44254606e-07f,
        1.02946107e-02f,
        5.87859135e-05f
    );

    float snimac2 = cubic(
        static_cast<float>(a2),
        -9.09109415e-10f,
        8.99002333e-07f,
        9.78867803e-03f,
        -9.65193864e-05f
    );

    ////////////////////////////////////////////////////////////////////
    // Control law area - update control variables (vent, spirala) based on snimac1, snimac2, globalTime, etc.

    vent = ...;
    spirala = ...;

    // 
    ///////////////////////////////////////////////////////////////////

    // End experiment after duration expires //////////////////////////

    // After experiment duration expires, force control inputs to zero
    bool experimentFinished = false;
    if (globalTime >= EXPERIMENT_DURATION_S) {
        vent = 0.0f;
        spirala = 0.0f;
        experimentFinished = true;
    }

    // Warning flag for too-large real period
    int warnTs = (timeTickPeriod > TS_WARNING_S) ? 1 : 0;

    // Print CSV line with all relevant data for logging and analysis ///

    // CSV output to Serial
    Serial.print(globalTime, 6);
    Serial.print(",");

    Serial.print(timeTickPeriod, 6);
    Serial.print(",");

    Serial.print(snimac1, 6);
    Serial.print(",");

    // Finish experiment and return to idle state
    if (experimentFinished) {
        stopOutputs();
        experimentRunning = false;
        Serial.println("DONE");
    }
}
\end{lstlisting}

\begin{lstlisting}[language=Python, caption={Externý skript pre logovanie dát}, label={lst:logger_py_appendix}]
import csv
import os
import sys
import time
from datetime import datetime

import serial


PORT = "COM3"          
BAUDRATE = 115200
TIMEOUT_S = 1.0
OUTPUT_DIR = "logs"

CSV_HEADER = "t,Ts,snimac1"


def make_output_path() -> str:
    os.makedirs(OUTPUT_DIR, exist_ok=True)
    timestamp = datetime.now().strftime("%Y-%m-%d_%H-%M-%S")
    return os.path.join(OUTPUT_DIR, f"experiment_{timestamp}.csv")


def print_red(text: str) -> None:
    print(f"\033[31m{text}\033[0m")


def main() -> None:
    output_path = make_output_path()

    print(f"Opening serial port: {PORT} @ {BAUDRATE}")
    print(f"Output file: {output_path}")

    try:
        with serial.Serial(PORT, BAUDRATE, timeout=TIMEOUT_S) as ser:
            # Let Arduino reset after opening the port
            time.sleep(2.0)

            print("Waiting for READY...")

            while True:
                raw_line = ser.readline()
                if not raw_line:
                    continue

                line = raw_line.decode("utf-8", errors="replace").strip()
                if not line:
                    continue

                print(f"RX: {line}")
                if line == "READY":
                    break

            with open(output_path, "w", newline="", encoding="utf-8") as csv_file:
                writer = csv.writer(csv_file)

                print("Sending START...")
                ser.write(b"START\n")
                ser.flush()

                header_received = False

                while True:
                    raw_line = ser.readline()
                    if not raw_line:
                        continue

                    line = raw_line.decode("utf-8", errors="replace").strip()
                    if not line:
                        continue

                    if line == "DONE":
                        print("Experiment finished.")
                        break

                    if line == "STOPPED":
                        print("Experiment stopped.")
                        break

                    if not header_received:
                        if line == CSV_HEADER:
                            writer.writerow(line.split(","))
                            csv_file.flush()
                            header_received = True
                            print("CSV header received. Logging started.")
                        else:
                            print(f"Skipping: {line}")
                        continue

                    parts = [part.strip() for part in line.split(",")]
                    writer.writerow(parts)
                    csv_file.flush()

                    warn_ts = len(parts) >= 11 and parts[10] == "1"
                    if warn_ts:
                        print_red(line)
                    else:
                        print(line)

    except serial.SerialException as exc:
        print(f"Serial error: {exc}")
        sys.exit(1)
    except KeyboardInterrupt:
        print("\nInterrupted by user.")


if __name__ == "__main__":
    main()
\end{lstlisting}